\chapter{高等代数}
    \begin{problem}
        假设 $n$ 维欧氏空间 $V$ 中有 $m$ 个向量两两成钝角, 
        证明: $m\leqslant n+1$.
    \end{problem}

    \begin{proof}
        当 $\mathrm{dim}V = 1$, 命题显然成立.

        假设当维数为 $n-1$ 时命题成立, 
        现证明命题在 $n$ 维时仍成立.
         设 $\ft{v}{1}{m}$ 为 $V$ 中两两成钝角的向量, 
        令 $W$ 为 $Lv_m$ 的正交补空间, 
        并记 $P:V\rightarrow W$ 为正交投影映射, 
        令 $w_1=Pv_1,\dots,w_{m-1}=Pv_{m-1}$, 则当 $i\neq j$ 时, 
        \begin{align*}
            \langle w_i, w_j \rangle 
            &= \langle Pv_i,Pv_j \rangle 
            = \left\langle v_i-\frac{\langle v_i,v_m \rangle}
            {\langle v_m,v_m \rangle}v_m, 
            v_j-\frac{\langle v_j,v_m \rangle}
            {\langle v_m,v_m \rangle}v_m\right\rangle \\
            &=\langle v_i,v_j \rangle - 
            \frac{\langle v_i,v_m \rangle\cdot
            \langle v_j,v_m \rangle}{\langle v_m,v_m \rangle} \\
            &<0
        \end{align*}
        因此 $\ft{w}{1}{m-1}$ 是 $n-1$ 维空间 $W$ 中两两成钝角的向量, 
        由归纳假设知 $m-1\leqslant n$, 即 $m\leqslant n+1$.
    \end{proof}
    \begin{remark}
        我感觉这和 $n$ 维欧氏空间的拓扑性质有关系, 
    \end{remark}
    \newpage

    \begin{problem}
        Let $V$ be an $n$-dimensional vector space over the field $F$, 
        then $\alpha\in\bigwedge^2(V^*)$ is decomposable, 
        i.e., there exist $\omega,\eta\in V^*$ 
        such that $\alpha=\omega\wedge\eta$, 
        if and only if $\alpha\wedge\alpha=0$.
    \end{problem}
    \begin{proof}
        Necessity is obvious. For sufficiency, we need a lemma:
        \begin{lemma}
            For $\alpha\in\bigwedge^r(V^*)$, $\omega\in V^*$,
            there exists $\eta\in\bigwedge^{r-1}(V^*)$ such that
            $\omega\wedge\eta=\alpha$ if and only if
            $\omega\wedge\alpha=0$.
        \end{lemma}
        If $\alpha\wedge\alpha=0$, 
        then 
        \[
            (\iota_X\alpha)\wedge\alpha 
            =\frac{1}{2}\iota_X(\alpha\wedge\alpha)=0,
            \quad\forall X\in V.
        \]
        Since $\alpha\neq0$, there exists $X\in V$ such that
        $\omega=\iota_X\alpha\neq0$. Then by the lemma, there exists
        $\eta\in\bigwedge^{1}(V^*)$ such that
        \[
            \omega\wedge\eta=\alpha.
        \]
    \end{proof}
    \begin{proof}[\indent Proof of lemma]
        Necessity is obvious.
        
        For sufficiency, let $\ft{\omega}{1}{n}$ be a basis of $V^*$
        with $\omega_1=\omega$.
        Then any $\alpha\in\bigwedge^r(V^*)$ can be written as 
        \[
            \alpha=\sum_{1\leqslant i_1<\cdots<i_r\leqslant n}
            a_{i_1\cdots i_r}\omega_{i_1}\wedge\cdots\wedge\omega_{i_r}.
        \]
        Then 
        \[
            \omega\wedge\alpha
            =\sum_{2\leqslant i_1<\cdots<i_r\leqslant n}
            a_{i_1\cdots i_r}\omega_1\wedge\omega_{i_1}\wedge\cdots\wedge\omega_{i_r}
            = 0,
        \]
        which implies that $a_{i_1\cdots i_r}=0$ for any
        $1<i_1<\cdots<i_r\leqslant n$. Thus
        \[
            \alpha
            =\sum_{2\leqslant i_2<\cdots<i_r\leqslant n}
            a_{1 i_2\cdots i_r}
            \omega_1\wedge\omega_{i_2}\wedge\cdots\wedge\omega_{i_r}
            = \omega\wedge\eta,
        \]
        where 
        \[
            \eta
            =\sum_{2\leqslant i_2<\cdots<i_r\leqslant n}
            a_{1 i_2\cdots i_r}
            \omega_{i_2}\wedge\cdots\wedge\omega_{i_r}.
        \]
        % \renewcommand{\qedsymbol}{}
    \end{proof}